%----------------------------------------------------------------------------------------
% GROEPSMORFISMEN. Alle opgaven zijn eigen werk.
%
% Opbouw: Probleem, Start, leerstof, en op een nieuwe pagina de Opdrachten, Koppel en Test jezelf.
% Labels van dit hoofdstuk beginnen met h3-.
%----------------------------------------------------------------------------------------

\mychapter{3}{Groepsmorfismen}

\begin{probleem}
Tot ongeveer 1975 rekenden ingenieurs met een rekenlat: twee latten met dezelfde schaal, waarop het getal $x$ staat op een afstand evenredig met $\log x$ van het begin. Schuif je de 1 van de bovenste lat boven de 2 van de onderste, dan staat onder elk getal $x$ van de bovenste lat het getal $2x$ op de onderste. In de figuur lees je onder de 3 het getal 6. Waarom werkt dat? En kun je een slimmere schaal bedenken waarmee een lat ook met negatieve getallen vermenigvuldigt, zodat je bijvoorbeeld $(-2) \cdot 3 = -6$ afleest?

\figuren{\begin{tikzpicture}[x=1cm,font=\scriptsize]
  \fill[accent!8] (-.25,-.75) rectangle (8.25,0);
  \draw[lijngrijs] (-.25,-.75) rectangle (8.25,0);
  \fill[accent!8] (2.408-.25,0) rectangle (2.408+8.25,.75);
  \draw[lijngrijs] (2.408-.25,0) rectangle (2.408+8.25,.75);
  \logschaal{0}{0}{8}{-1}
  \logschaal{2.408}{0}{8}{1}
  \draw[fel!85!black,thick] (2.408,-.85) -- (2.408,.85);
  \draw[fel!85!black,thick] (6.225,-.85) -- (6.225,.85);
  \node[anchor=west,text=grijs] at (10.75,.4) {bovenste lat};
  \node[anchor=west,text=grijs] at (8.35,-.4) {onderste lat};
\end{tikzpicture}}
\end{probleem}

\lkr{Leerplan Complementair (III-WiCo-d): LPD 32, het deel over groepsmorfismen, met de wenk om isomorfismen te behandelen en voorbeelden te geven van isomorfe groepen. Het hoofdstuk definieert ook morfismen in het algemeen (sectie 1), omdat ze in pakket 4 al voorkomen zonder naam: de logaritme, de determinant, de rest bij deling. De kern, het product van groepen en het overzicht van de kleine groepen zijn uitbreidingen. Voorkennis: hoofdstuk 1 en 2, de rekenregels van logaritmen (Exponentiële en logaritmische functies) en het idee ``dezelfde structuur'' uit Groepen, hoofdstuk 3.}

In Groepen zag je dat twee groepen dezelfde structuur kunnen hebben: geef je de elementen van de ene groep een andere naam, dan krijg je de tabel van de andere. Zo is een draaiing van de gelijkzijdige driehoek ``hetzelfde'' als een permutatie van haar hoekpunten. Je zag ook dat de logaritme een product omzet in een som. In beide gevallen gaat een functie van de ene groep naar de andere, en ze respecteert de bewerking.

Zo'n functie heet een \textit{groepsmorfisme}. In dit hoofdstuk leer je herkennen wanneer een functie een morfisme is, en wat ze dan met het neutraal element, de inversen en de orde van een element doet. Is het morfisme bovendien een bijectie, dan is het een \textit{isomorfisme}: de twee groepen zijn dan dezelfde groep in een andere vermomming. Met wat een isomorfisme bewaart, toon je aan dat twee groepen niet isomorf zijn. Zo vind je ook waarom een rekenlat niet met negatieve getallen kan werken. Op het einde zie je dat er maar weinig verschillende kleine groepen bestaan.

\begin{oefeningen}[Start]
\oef{h3-start-1}

\oef{h3-verken}
\end{oefeningen}

\lkr{Laat oefening \ref{oef:h3-verken} maken vóór de uitleg, met een geodriehoek of een lat. De leerlingen tekenen zelf een logaritmische schaal en zien dat de afstanden van 2 en 4 samen de afstand van 8 geven: de logaritme zet vermenigvuldigen om in optellen.}

%========================================================================================
\section{Groepsmorfismen}

\subsection{Functies die de bewerking respecteren}

In de Verken telde je op je eigen lat de afstanden van 2 en 4 op, en kwam je uit bij 8. Waarom komt er een product uit? Het getal $x$ staat op afstand $10 \log x$ van het begin. Leg je twee afstanden achter elkaar, dan tel je ze op: $10 \log a + 10 \log b = 10 \log(ab)$. Je kunt ook in de omgekeerde richting rekenen, want elke afstand hoort bij precies één getal. Dat de logaritme van een product de som is van de logaritmen, zegt dat de functie $\log$ de bewerking van de ene groep $(\R^+_0, \cdot)$ omzet in die van de andere groep $(\R, +)$.

\begin{theorie}
Neem twee groepen $(G, \ast)$ en $(H, \circ)$. Een functie $f \colon G \to H$ is een \textbf{\textcolor{accent}{groepsmorfisme}} als

\vspace{1mm}
\qquad $\formule{f(a \ast b) = f(a) \circ f(b)}$ \qquad voor alle $a, b \in G$.
\end{theorie}

In woorden: het maakt niet uit of je eerst in $G$ rekent en dan $f$ toepast, of eerst $f$ toepast en dan in $H$ rekent. Je vindt het woord ook als \textit{homomorfisme}. Let op de bewerkingen: links staat de bewerking van $G$, rechts die van $H$. Bij de logaritme is dat
\[ \log(a \cdot b) = \log a + \log b . \]

\figuren{\begin{tikzpicture}[font=\small]
  \node (a) at (0,1.4) {$a, \ b$};
  \node (ab) at (4.2,1.4) {$a \ast b$};
  \node (fa) at (0,0) {$f(a), \ f(b)$};
  \node (fab) at (4.2,0) {$f(a \ast b) = f(a) \circ f(b)$};
  \draw[-{Stealth[length=2mm]},accent,thick] (a) -- node[above] {rekenen in $G$} (ab);
  \draw[-{Stealth[length=2mm]},accent,thick] (fa) -- node[below] {rekenen in $H$} (fab);
  \draw[-{Stealth[length=2mm]},fel!85!black,thick] (a) -- node[left] {$f$} (fa);
  \draw[-{Stealth[length=2mm]},fel!85!black,thick] (ab) -- node[right] {$f$} (fab);
\end{tikzpicture}}

Je kent al veel morfismen, zonder dat ze die naam droegen. Zo zegt de rekenregel $a^m \ast a^n = a^{m + n}$ uit hoofdstuk 1 dat $k \mapsto a^k$ een morfisme is van $(\Z, +)$ naar $(G, \ast)$, en de regel $\det(AB) = \det A \cdot \det B$ dat de determinant een morfisme is van de inverteerbare matrices naar $(\R_0, \cdot)$.

\voorbeeld{}{Toon aan dat $f \colon (\Z, +) \to (\Zn{12}, \plusn)$, met $f(k)$ de rest van $k$ bij deling door 12, een groepsmorfisme is.\label{vb:h3-rest}}{%
\uitleg{Links tellen we gewone gehele getallen op, rechts rekenen we modulo 12. We schrijven $a$ en $b$ met hun rest, en tonen dat $f(a + b) = f(a) \plusn f(b)$.}
\begin{stappenbreed}
a = 12q + f(a), \quad b = 12q' + f(b) & Deling met rest. \\[1ex]
a + b = 12(q + q') + f(a) + f(b) & Dus $a + b$ en $f(a) + f(b)$ verschillen een veelvoud van 12. \\[1ex]
f(a + b) = \text{rest van } f(a) + f(b) = f(a) \plusn f(b) & Getallen die een veelvoud van 12 verschillen, hebben dezelfde rest.
\end{stappenbreed}
Dit is precies wat je in Groepen deed: eerst gewoon rekenen, en pas op het einde de rest nemen. Zo is $f(17 + 20) = f(37) = 1$ en $f(17) \plusn f(20) = 5 \plusn 8 = 1$.}

\voorbeeld{}{Is $f \colon (\R, +) \to (\R^+_0, \cdot)$ met $f(x) = 2^x$ een groepsmorfisme? En $g \colon (\R, +) \to (\R, +)$ met $g(x) = x^2$?\label{vb:h3-morfisme}}{%
\uitleg{We vergelijken $f(a + b)$ met $f(a) \cdot f(b)$: links de bewerking van de eerste groep, rechts die van de tweede. Bij $g$ is de bewerking twee keer een optelling.}
\begin{stappen}
f(a + b) = 2^{a + b} = 2^a \cdot 2^b = f(a) \cdot f(b) & Een rekenregel van machten: $f$ is een morfisme. \\[1ex]
g(1 + 1) = 4, \ \ g(1) + g(1) = 2 & Eén tegenvoorbeeld: $g$ is geen morfisme.
\end{stappen}
De exponentiële functie zet een som om in een product, de logaritme doet het omgekeerde.}

\begin{werkwijze}
Schrijf $f(a \ast b)$ en $f(a) \circ f(b)$ uit met de bewerking van elke groep, en toon dat ze gelijk zijn voor alle $a$ en $b$. Om te tonen dat $f$ geen morfisme is, volstaat één koppel $(a, b)$.
\end{werkwijze}

\begin{oefeningen}
\oef{h3-morfisme}
\end{oefeningen}

\subsection{Wat een morfisme doet met $e$ en met inversen}

Bij de logaritme weet je dat $\log 1 = 0$ en $\log \frac1a = -\log a$: het neutraal element gaat naar het neutraal element, een inverse naar de inverse. Is dat toeval, of geldt het voor elk morfisme?

\begin{theorie}
\begin{stelling}{Beeld van het neutraal element}
Een groepsmorfisme $f \colon G \to H$ beeldt het neutraal element $e$ van $G$ af op het neutraal element $e'$ van $H$: \qquad $\formule{f(e) = e'}$
\end{stelling}

\begin{bewijs}
Uit $e \ast e = e$ volgt
\bewijsregel{f(e) \circ f(e) = f(e \ast e) = f(e) = f(e) \circ e' ,}
en schrappen in $H$ geeft $f(e) = e'$.
\end{bewijs}
\end{theorie}

\begin{theorie}
\begin{stelling}{Beeld van een inverse en van een macht}
Voor een groepsmorfisme $f$ is \qquad $\formule{f(a^{-1}) = f(a)^{-1}}$ \qquad en \qquad $\formule{f(a^k) = f(a)^k}$ \quad voor elk geheel getal $k$.
\end{stelling}

\begin{bewijs}
$f(a) \circ f(a^{-1}) = f(a \ast a^{-1}) = f(e) = e'$, en net zo $f(a^{-1}) \circ f(a) = e'$. Dus $f(a^{-1})$ is de inverse van $f(a)$. Voor $k \geqslant 0$ volgt $f(a^k) = f(a)^k$ door $k$ keer de definitie toe te passen, en voor $k < 0$ is $f(a^k) = f\bigl((a^{-1})^{-k}\bigr) = \bigl(f(a)^{-1}\bigr)^{-k} = f(a)^k$.
\end{bewijs}
\end{theorie}

Een morfisme ligt dus vast door weinig gegevens. Ken je het beeld van een voortbrenger, dan ken je het beeld van elke macht: $f(a^k) = f(a)^k$. Voor de symmetriegroep van de gelijkzijdige driehoek schrijven we vanaf nu kort $D_3$: in het algemeen is $D_n$ de symmetriegroep van de regelmatige $n$-hoek, met $2n$ elementen.

\voorbeeld{}{$f \colon (\Z, +) \to D_3$ is een morfisme naar de symmetriegroep van de gelijkzijdige driehoek, met $f(1) = r$. Bereken $f(0)$, $f(-1)$ en $f(5)$. Welke getallen hebben als beeld $e$?\label{vb:h3-beeld}}{%
\uitleg{In $(\Z, +)$ is een ``macht'' van 1 een veelvoud: $5 = 5 \cdot 1$ en $-1$ is de inverse van 1. We passen de stellingen toe.}
\begin{stappen}
f(0) = e & Beeld van het neutraal element. \\[1ex]
f(-1) = f(1)^{-1} = r^{-1} = r^2 & Beeld van een inverse. \\[1ex]
f(5) = f(1)^5 = r^5 = r^2 & $r^3 = e$. \\[1ex]
f(k) = r^k = e \iff 3 \mid k & Wanneer is $a^k = e$? (hoofdstuk 1)
\end{stappen}
De veelvouden van 3 hebben als beeld $e$. Ze vormen de deelgroep $3\Z$.}


\begin{oefeningen}
\oef{h3-beeld}
\end{oefeningen}

\subsection{De kern}

In het voorbeeld vormden de elementen met beeld $e$ een deelgroep. Dat is altijd zo, en die deelgroep vertelt hoeveel informatie het morfisme ``vergeet''.

\begin{theorie}
De \textbf{\textcolor{accent}{kern}} van een groepsmorfisme $f \colon G \to H$ is \qquad $\formule{\ker f = \{a \in G \mid f(a) = e'\}}$

\begin{stelling}{De kern is een deelgroep}
$\ker f$ is een deelgroep van $G$.
\end{stelling}

\begin{bewijs}
$e \in \ker f$, want $f(e) = e'$. Liggen $a$ en $b$ in de kern, dan is $f(a \ast b) = e' \circ e' = e'$ en $f(a^{-1}) = (e')^{-1} = e'$. Volgens het kenmerk van een deelgroep is $\ker f$ een deelgroep.
\end{bewijs}
\end{theorie}

\begin{theorie}
\begin{stelling}{Injectief morfisme}
Een groepsmorfisme $f$ is injectief \qquad $\formule{\iff \ker f = \{e\}}$
\end{stelling}

\begin{bewijs}
Is $f$ injectief, dan heeft alleen $e$ het beeld $e'$. Omgekeerd, stel $\ker f = \{e\}$ en $f(a) = f(b)$. Dan is
\bewijsregel{f(a \ast b^{-1}) = f(a) \circ f(b)^{-1} = e' ,}
dus $a \ast b^{-1} \in \ker f$, zodat $a \ast b^{-1} = e$ en $a = b$.
\end{bewijs}
\end{theorie}

Bij een morfisme volstaat het dus na te gaan welke elementen op $e'$ terechtkomen. Is dat alleen $e$, dan gaan verschillende elementen nooit naar hetzelfde beeld.

\voorbeeld{}{Bepaal de kern van $f \colon (\Zn{12}, \plusn) \to (\Zn{4}, \plusn)$, met $f(k)$ de rest van $k$ bij deling door 4. Is $f$ injectief?\label{vb:h3-kern}}{%
\uitleg{Omdat 4 een deler is van 12, is $f$ een morfisme, zoals in voorbeeld \ref{vb:h3-rest}. In de kern zitten de elementen van $\Zn{12}$ met rest 0 bij deling door 4.}
\begin{stappen}
\ker f = \{0, 4, 8\} & De veelvouden van 4 in $\Zn{12}$. \\[1ex]
\ker f \neq \{0\} & $f$ is niet injectief: zo is $f(1) = f(5) = 1$.
\end{stappen}
De kern is de deelgroep $\vb{4}$ met 3 elementen, en elk element van $\Zn{4}$ is het beeld van precies 3 elementen: de nevenklasse $k \plusn \ker f$.}

\begin{oefeningen}
\oef{h3-kern}
\end{oefeningen}

%========================================================================================
\section{Isomorfismen}

\subsection{Dezelfde groep in een andere vermomming}

In Groepen had $(\Zn{4}, \plusn)$ dezelfde structuur als $\bigl(\{1, i, -1, -i\}, \cdot\bigr)$: geef $k$ de naam $i^k$, en de tabel van de ene wordt de tabel van de andere. Die nieuwe naam is een functie die elk element op precies één element afbeeldt, en die de bewerking respecteert.

\begin{theorie}
Een \textbf{\textcolor{accent}{isomorfisme}} is een groepsmorfisme dat bijectief is.

\vspace{1mm}
Bestaat er een isomorfisme van $G$ naar $H$, dan zijn $G$ en $H$ \textbf{\textcolor{accent}{isomorf}}: \quad $\formule{G \isom H}$

\begin{stelling}{Inverse van een isomorfisme}
Is $f \colon G \to H$ een isomorfisme, dan is ook $f^{-1} \colon H \to G$ een isomorfisme.
\end{stelling}

\begin{bewijs}
$f^{-1}$ is een bijectie. Neem $x, y \in H$, en schrijf $x = f(a)$ en $y = f(b)$. Dan is $f(a \ast b) = x \circ y$, dus
\bewijsregel{f^{-1}(x \circ y) = a \ast b = f^{-1}(x) \ast f^{-1}(y) .}
\end{bewijs}
\end{theorie}

Isomorf is dus een verband in beide richtingen: is $G \isom H$, dan ook $H \isom G$. Is een groep isomorf met een tweede, en die met een derde, dan is de eerste ook isomorf met de derde: stel de twee isomorfismen samen. Voor een wiskundige zijn isomorfe groepen ``dezelfde'' groep. Alles wat je met de bewerking kunt zeggen, zoals de orde van een element of welke elementen wisselen, geldt in de ene als en slechts als het in de andere geldt.

\voorbeeld{}{Toon aan dat $f \colon (\Zn{4}, \plusn) \to \bigl(\{1, i, -1, -i\}, \cdot\bigr)$ met $f(k) = i^k$ een isomorfisme is.\label{vb:h3-iso}}{%
\uitleg{We tonen dat $f$ een bijectie is en dat $f(k \plusn l) = f(k) \cdot f(l)$. De moeilijkheid: links nemen we de rest bij deling door 4, rechts niet.}
\begin{stappenbreed}
f(0) = 1, \ f(1) = i, \ f(2) = -1, \ f(3) = -i & Vier verschillende beelden: een bijectie. \\[1ex]
k + l = 4q + (k \plusn l) & Deling met rest. \\[1ex]
f(k) \cdot f(l) = i^{k + l} = (i^4)^q \cdot i^{k \plusn l} = f(k \plusn l) & Want $i^4 = 1$.
\end{stappenbreed}
Dus $(\Zn{4}, \plusn) \isom \bigl(\{1, i, -1, -i\}, \cdot\bigr)$.}

\begin{werkwijze}
Om te tonen dat $G \isom H$, zoek je een functie $f \colon G \to H$, en toon je dat ze een morfisme is en bijectief. Voor eindige groepen met evenveel elementen volstaat injectief: kijk of $\ker f = \{e\}$. Een goede kandidaat stuurt een voortbrenger naar een voortbrenger, of een element naar een element met dezelfde orde.
\end{werkwijze}

\begin{oefeningen}
\oef{h3-iso}
\end{oefeningen}

\subsection{Wat een isomorfisme bewaart}

Hoe toon je aan dat twee groepen \textit{niet} isomorf zijn? Je kunt niet alle functies tussen twee groepen nagaan. Maar een isomorfisme bewaart alles wat met de bewerking te maken heeft. Vind je een eigenschap die de ene groep heeft en de andere niet, dan zijn ze niet isomorf.

\begin{theorie}
\begin{stelling}{Isomorfisme en orde}
Is $f \colon G \to H$ een isomorfisme, dan is \qquad $\formule{\ord\bigl(f(a)\bigr) = \ord(a)}$ \qquad voor elk element $a$.
\end{stelling}

\begin{bewijs}
$f(a)^k = f(a^k)$, en dat is $e'$ als en slechts als $a^k \in \ker f = \{e\}$. Dus $f(a)^k = e' \iff a^k = e$: dezelfde exponenten $k$, dus dezelfde kleinste. Ook een oneindige orde blijft oneindig.
\end{bewijs}
\end{theorie}

\begin{theorie}
\begin{stelling}{Isomorfisme en commutatief}
Is $G \isom H$ en is $G$ commutatief, dan is ook $H$ commutatief.
\end{stelling}

\begin{bewijs}
Neem $x = f(a)$ en $y = f(b)$ in $H$. Dan is $x \circ y = f(a \ast b) = f(b \ast a) = y \circ x$.
\end{bewijs}
\end{theorie}

Isomorfe groepen hebben dus evenveel elementen, evenveel elementen van elke orde, en ze zijn allebei commutatief of allebei niet. Dat is wat je in Groepen al gebruikte om groepen te vergelijken.

\voorbeeld{}{Zijn $(\Zn{4}, \plusn)$ en $\bigl(\Zs{8}, \maaln\bigr) = \bigl(\{1, 3, 5, 7\}, \maaln\bigr)$ isomorf? En $(\Zn{6}, \plusn)$ en $S_3$?\label{vb:h3-nietiso}}{%
\uitleg{De groepen hebben telkens evenveel elementen. We zoeken een eigenschap die door een isomorfisme bewaard wordt, maar niet in beide groepen geldt.}
\begin{stappenbreed}
\ord(1) = 4 \text{ in } \Zn{4}; \ \ 3^2 = 5^2 = 7^2 = 1 \text{ in } \Zs{8} & $\Zn{4}$ heeft een element van orde 4, $\Zs{8}$ niet: niet isomorf. \\[1ex]
\cyc{1 & 2} \circ \cyc{1 & 3} \neq \cyc{1 & 3} \circ \cyc{1 & 2} & $S_3$ is niet commutatief, $\Zn{6}$ wel: niet isomorf.
\end{stappenbreed}}

\begin{werkwijze}
Om te tonen dat $G$ en $H$ niet isomorf zijn, zoek je een eigenschap die een isomorfisme bewaart en die maar in één van beide geldt: het aantal elementen, commutatief of niet, cyclisch of niet, het aantal elementen van een gegeven orde.
\end{werkwijze}

\begin{oefeningen}
\oef{h3-nietiso}
\end{oefeningen}

\voorbeeld{}{Terug naar het probleem. Tot ongeveer 1975 rekenden ingenieurs met een rekenlat: twee latten met dezelfde schaal, waarop het getal $x$ staat op een afstand evenredig met $\log x$ van het begin. Schuif je de 1 van de bovenste lat boven de 2 van de onderste, dan staat onder elk getal $x$ van de bovenste lat het getal $2x$ op de onderste. In de figuur lees je onder de 3 het getal 6. Waarom werkt dat? En kun je een slimmere schaal bedenken waarmee een lat ook met negatieve getallen vermenigvuldigt, zodat je bijvoorbeeld $(-2) \cdot 3 = -6$ afleest?

\figuren{\begin{tikzpicture}[x=1cm,font=\scriptsize]
  \fill[accent!8] (-.25,-.75) rectangle (8.25,0);
  \draw[lijngrijs] (-.25,-.75) rectangle (8.25,0);
  \fill[accent!8] (2.408-.25,0) rectangle (2.408+8.25,.75);
  \draw[lijngrijs] (2.408-.25,0) rectangle (2.408+8.25,.75);
  \logschaal{0}{0}{8}{-1}
  \logschaal{2.408}{0}{8}{1}
  \draw[fel!85!black,thick] (2.408,-.85) -- (2.408,.85);
  \draw[fel!85!black,thick] (6.225,-.85) -- (6.225,.85);
  \node[anchor=west,text=grijs] at (10.75,.4) {bovenste lat};
  \node[anchor=west,text=grijs] at (8.35,-.4) {onderste lat};
\end{tikzpicture}}\label{vb:h3-terug}}{%
\uitleg{Een lat met een schaal is een functie die elk getal op een plaats zet, en verschillende getallen op verschillende plaatsen: een injectieve functie $f$ die het getal $x$ op afstand $f(x)$ van het begin zet. De lat vermenigvuldigt als $f(a \cdot b) = f(a) + f(b)$: dan is $f$ een injectief morfisme. Voor de negatieve getallen moet $f$ werken op de groep $(\R_0, \cdot)$, en we kijken wat het morfisme doet met $-1$.}
\begin{stappenbreed}
f(x) = L \log x & De logaritmische schaal, met $L$ de lengte van de lat. \\[1ex]
f(2 \cdot 3) = L \log 2 + L \log 3 = f(2) + f(3) & Schuif de 1 van de bovenste lat (afstand 0) boven de 2: onder $x$ staat het getal op afstand $f(2) + f(x)$, dus $2x$. \\[1ex]
(-1) \cdot (-1) = 1 & $-1$ heeft orde 2 in $(\R_0, \cdot)$. \\[1ex]
f(-1) + f(-1) = f(1) = 0 & $f$ is een morfisme naar $(\R, +)$, en $f(1)$ is het neutraal element 0. \\[1ex]
f(-1) = 0 = f(1) & In $(\R, +)$ is $x + x = 0$ alleen voor $x = 0$: geen element van orde 2. \\[1ex]
-1 \neq 1 & Maar $f$ moest injectief zijn. Tegenspraak.
\end{stappenbreed}
De rekenlat werkt omdat $\log$ een isomorfisme is van $(\R^+_0, \cdot)$ naar $(\R, +)$: elk getal heeft een plaats, en elke plaats hoort bij een getal. Een schaal voor alle reële getallen verschillend van 0 bestaat niet, hoe slim ook: $(\R_0, \cdot)$ heeft een element van orde 2 en $(\R, +)$ niet. Daarom rekent wie een rekenlat gebruikt het teken apart uit: min maal min is plus.}

%========================================================================================
\section{Kleine groepen}

\subsection{Cyclische groepen}

Hoeveel ``verschillende'' groepen zijn er met een gegeven aantal elementen? Voor een wiskundige tellen isomorfe groepen als dezelfde. Bij de cyclische groepen is het antwoord eenvoudig: alle cyclische groepen met $n$ elementen zijn isomorf.

\begin{theorie}
\begin{stelling}{Cyclische groepen}
Elke cyclische groep met $n$ elementen is isomorf met $(\Zn{n}, \plusn)$. Elke oneindige cyclische groep is isomorf met $(\Z, +)$.
\end{stelling}

\begin{bewijs}
Stel $G = \vb{a}$ met $\ord(a) = n$, en neem $f \colon \Zn{n} \to G$ met $f(k) = a^k$. De $n$ machten $e, a, \ldots, a^{n - 1}$ zijn verschillend (hoofdstuk 1), dus $f$ is bijectief. Omdat $k + l$ en $k \plusn l$ een veelvoud van $n$ verschillen, is (gelijke machten)
\bewijsregel{f(k) \ast f(l) = a^{k + l} = a^{k \plusn l} = f(k \plusn l) .}
Heeft $a$ een oneindige orde, neem dan $f \colon \Z \to G$ met $f(k) = a^k$: alle machten zijn verschillend, en $a^k \ast a^l = a^{k + l}$.
\end{bewijs}
\end{theorie}

Met de stelling van Lagrange volgt: elke groep met een priemgetal $p$ als orde is cyclisch, dus isomorf met $(\Zn{p}, \plusn)$. Er is maar één groep met 2, 3, 5 of 7 elementen. En een isomorfisme vind je door een voortbrenger naar een voortbrenger te sturen.

\voorbeeld{}{Geef een isomorfisme van $(\Zn{6}, \plusn)$ naar $\bigl(\Zs{7}, \maaln\bigr)$.\label{vb:h3-cyclisch}}{%
\uitleg{In hoofdstuk 1 vonden we dat 3 een voortbrenger is van $\Zs{7} = \{1, 2, \ldots, 6\}$. Volgens de stelling is $f(k) = 3^k$ een isomorfisme.}
\begin{stappen}
f(0), \ldots, f(5) = 1, \ 3, \ 2, \ 6, \ 4, \ 5 & De machten van 3 modulo 7. \\[1ex]
f(2 \plusn 5) = f(1) = 3 & Een controle. \\[1ex]
f(2) \maaln f(5) = 2 \maaln 5 = 10 \to 3 & Hetzelfde.
\end{stappen}
Met de voortbrenger 5 krijg je een ander isomorfisme: er zijn er evenveel als voortbrengers.}

\begin{oefeningen}
\oef{h3-cyclisch}
\end{oefeningen}

\subsection{Het product van twee groepen}

Uit twee groepen kun je een nieuwe groep maken door koppels te nemen en in elke plaats apart te rekenen. Zo kun je groepen bouwen die niet cyclisch zijn.

\begin{theorie}
Het \textbf{\textcolor{accent}{product}} van de groepen $(G, \ast)$ en $(H, \circ)$ is de verzameling van de koppels $(g, h)$ met $g \in G$ en $h \in H$, met de bewerking

\vspace{1mm}
\qquad $\formule{(g, h) \cdot (g', h') = (g \ast g', \ h \circ h')}$

\vspace{1mm}
We noteren het product $\formule{G \times H}$. Het is een groep met $|G| \cdot |H|$ elementen.
\end{theorie}

De voorwaarden van een groep gelden in elke plaats apart, dus ook voor de koppels: het neutraal element is $(e, e')$ en de inverse van $(g, h)$ is $(g^{-1}, h^{-1})$. Het kleinste voorbeeld dat niet cyclisch is, is $\Zn{2} \times \Zn{2}$: elk koppel opgeteld met zichzelf geeft $(0, 0)$.

\voorbeeld{}{Toon aan dat $\Zn{2} \times \Zn{3}$ cyclisch is en dat $\Zn{2} \times \Zn{2}$ isomorf is met de symmetriegroep van de rechthoek.\label{vb:h3-product}}{%
\uitleg{Voor de eerste groep zoeken we een element van orde 6. Voor de tweede geven we elk koppel de naam van een symmetrie: $h$ en $v$ krijgen de twee koppels met één 1, de halve draai $d = h \circ v$ hun som.}
\begin{stappenbreed}
(1, 1), \ (0, 2), \ (1, 0), \ (0, 1), \ (1, 2), \ (0, 0) & De veelvouden van $(1, 1)$ in $\Zn{2} \times \Zn{3}$: orde 6, dus cyclisch en isomorf met $\Zn{6}$. \\[1ex]
(0, 0) \mapsto e, \ (1, 0) \mapsto h, \ (0, 1) \mapsto v, \ (1, 1) \mapsto d & De nieuwe namen.
\end{stappenbreed}
In beide groepen is elk element zijn eigen inverse, en geven twee verschillende elementen die niet het neutraal element zijn samen het derde: $(1, 0) + (0, 1) = (1, 1)$ en $h \circ v = d$. Beide tabellen hebben dus dezelfde vorm: $\Zn{2} \times \Zn{2}$ is isomorf met de symmetriegroep van de rechthoek. Men noemt deze groep ook de \textit{viergroep van Klein}.}

\begin{oefeningen}
\oef{h3-product}
\end{oefeningen}

\subsection{Alle kleine groepen}

Met de stellingen van dit hoofdstuk kun je bepalen hoeveel niet-isomorfe groepen er zijn met een klein aantal elementen. Voor een priemgetal is er maar één. Voor 4 elementen zijn er twee: $\Zn{4}$ en de viergroep van Klein (Groepen, hoofdstuk 3). Voor 6 elementen ook twee: $\Zn{6}$ en $S_3$ (zie de opdrachten).

\begin{center}
\renewcommand{\arraystretch}{1.3}
\begin{tabular}{|c|c|c|}
\hline
\rowcolor{accent!14} orde & groepen, op isomorfisme na & voorbeelden uit deze cursus \\ \hline
1, 2, 3, 5, 7 & $\Zn{n}$ & de draaiingen van een regelmatige $n$-hoek \\ \hline
4 & $\Zn{4}$, \ $\Zn{2} \times \Zn{2}$ & $\{1, i, -1, -i\}$; de rechthoek, $\Zs{8}$, $\Zs{12}$ \\ \hline
6 & $\Zn{6}$, \ $S_3$ & $\Zs{7}$, $\Zs{9}$; de gelijkzijdige driehoek \\ \hline
8 & $\Zn{8}$, \ $\Zn{4} \times \Zn{2}$, \ $\Zn{2} \times \Zn{2} \times \Zn{2}$, \ $D_4$, \ $Q_8$ & $\Zs{15} \isom \Zn{4} \times \Zn{2}$; het vierkant \\ \hline
\end{tabular}
\end{center}

Met 8 elementen zijn er vijf groepen: drie commutatieve en twee die niet commutatief zijn. De symmetriegroep $D_4$ van het vierkant ken je. De laatste, de \textit{quaterniongroep} $Q_8$, bestaat uit $\pm 1$, $\pm i$, $\pm j$ en $\pm k$ met $i^2 = j^2 = k^2 = ijk = -1$: getallen die William Rowan Hamilton in 1843 bedacht en in een brug in Dublin kerfde. Dat er niet meer zijn, bewijzen we niet. Voor grotere aantallen groeit het snel: er zijn 14 groepen met 16 elementen, en meer dan 49 miljard met 1024 elementen.

%========================================================================================
\opdrachtenkop

\begin{oefeningen}
\oef{h3-o1}

\oef{h3-o2}

\oef{h3-o3}

\oef{h3-o4}

\oef{h3-o5}

\oef{h3-o6}

\oef{h3-o7}

\oef{h3-o8}

\oef{h3-o9}

\oef{h3-o10}

\oef{h3-o11}

\oef{h3-o12}

\oef{h3-o13}

\oef{h3-o14}

\oef{h3-o15}

\oef{h3-o16}

\oef{h3-o17}

\oef{h3-o18}

\oef{h3-o19}

\oef{h3-o20}

\oef{h3-o21}

\oef{h3-o22}
\end{oefeningen}

%----------------------------------------------------------------------------------------
% KOPPEL: verbanden met andere leerstof
%----------------------------------------------------------------------------------------
\needspace{20\baselineskip}
\opdrachtsectie{Koppel}

\lkr{Exponentiële en logaritmische functies, Complexe getallen en Matrices zitten in het 5de jaar van pakket 4.}

\begin{oefeningen}
\oef{h3-koppel-1}

\oef{h3-koppel-2}

\oef{h3-koppel-3}
\end{oefeningen}

%----------------------------------------------------------------------------------------
\needspace{30\baselineskip}
\opdrachtsectie{Test jezelf}

\begin{oefeningen}
\oef{h3-t1}

\oef{h3-t2}

\oef{h3-t3}

\oef{h3-t4}

\oef{h3-t5}

\oef{h3-t6}

\oef{h3-t7}
\end{oefeningen}

\herhaalschema{%
  \fv{oef:h3-t1}{vb:h3-morfisme}{oef:h3-morfisme}%
  \fv{oef:h3-t2}{vb:h3-beeld}{oef:h3-beeld}%
  \fv{oef:h3-t3}{vb:h3-kern}{oef:h3-kern}%
  \fv{oef:h3-t4}{vb:h3-iso}{oef:h3-iso}%
  \fv{oef:h3-t5}{vb:h3-nietiso}{oef:h3-nietiso}%
  \fv{oef:h3-t6}{vb:h3-beeld}{oef:h3-o2}%
  \fv{oef:h3-t7}{vb:h3-cyclisch}{oef:h3-cyclisch}}

\ifzelfstudie
\section*{Maak je eigen samenvatting}
Sluit dit boek en schrijf in je schrift: wat een groepsmorfisme is, met drie voorbeelden en een tegenvoorbeeld, wat een morfisme doet met het neutraal element, met inversen en met machten, wat de kern is en wanneer een morfisme injectief is, wat een isomorfisme is, wat een isomorfisme bewaart en hoe je toont dat twee groepen niet isomorf zijn, waarom een rekenlat niet met negatieve getallen werkt, waarom alle cyclische groepen met $n$ elementen isomorf zijn met $\Zn{n}$, wat het product van twee groepen is, en welke groepen er zijn met 4 en met 6 elementen. Vergelijk daarna met de kaders.
\fi
